\section*{Chapitre \ref{ch:g10:chemistry-of-life} --- La composition chimique des êtres vivants}

\begin{solution}{exo:g10:chemistry-of-life:1}
Oxygène (environ 65~\%), carbone (18.5~\%), hydrogène (9.5~\%) et azote
(3.2~\%)~: 96~\% de la masse du corps ensemble.
\end{solution}

\begin{solution}{exo:g10:chemistry-of-life:2}
Minéral~: eau, chlorure de sodium, phosphate de calcium, dioxyde de
carbone (un composé du carbone, mais produit hors du vivant et sans
squelette carbone--hydrogène). Organique~: glucose, triglycéride,
\omterm{def:g10:chemistry-of-life:families}{acide aminé}.
\end{solution}

\begin{solution}{exo:g10:chemistry-of-life:3}
Eau perdue $12.0 - 0.6 = \qty{11.4}{g}$~; $11.4/12.0 = 0.95$~: la
feuille est à 95~\% d'eau.
\end{solution}

\begin{solution}{exo:g10:chemistry-of-life:4}
(a) Eau iodée sur la graine écrasée~: bleu-noir. (b) Réactif du
biuret~: violet. (c) Écraser la graine sur du papier~: une tache
translucide qui ne sèche pas~; le Soudan III la teinte en rouge.
\end{solution}

\begin{solution}{exo:g10:chemistry-of-life:5}
\omterm{def:g10:chemistry-of-life:families}{Glucides}~: sucres simples comme le glucose. \omterm{def:g10:chemistry-of-life:families}{Lipides}~: glycérol et
acides gras. \omterm{def:g10:chemistry-of-life:families}{Protéines}~: \omterm{def:g10:chemistry-of-life:families}{acides aminés} (vingt sortes). \omterm{def:g10:chemistry-of-life:families}{Acides
nucléiques}~: nucléotides (quatre sortes).
\end{solution}

\begin{solution}{exo:g10:chemistry-of-life:6}
Carbone~: $18.5/0.02 \approx 900$ fois plus concentré. Azote~:
$3.2/0.002 = 1600$ fois. La matière vivante n'est pas un échantillon de
son environnement~: elle rassemble quelques éléments rares d'un facteur
mille, ce qui est une signature chimique de la vie.
\end{solution}

\begin{solution}{exo:g10:chemistry-of-life:7}
Eau~: $0.62 \times 70 = \qty{43.4}{kg}$. Matière sèche~: $70 - 43.4 =
\qty{26.6}{kg}$. \omterm{def:g10:chemistry-of-life:families}{Protéines}~: $0.17 \times 70 = \qty{11.9}{kg}$.
\end{solution}

\begin{solution}{exo:g10:chemistry-of-life:8}
\omterm{def:g10:chemistry-of-life:families}{Glucides} $60 \times 17 = \qty{1020}{kJ}$~; \omterm{def:g10:chemistry-of-life:families}{lipides} $12 \times 38 =
\qty{456}{kJ}$~; \omterm{def:g10:chemistry-of-life:families}{protéines} $8 \times 17 = \qty{136}{kJ}$~; total
\qty{1612}{kJ}, soit environ \qty{1600}{kJ}. Les \omterm{def:g10:chemistry-of-life:families}{lipides} fournissent $456/1612 \approx
28\%$ de l'énergie avec seulement 12~\% de la masse.
\end{solution}

\begin{solution}{exo:g10:chemistry-of-life:9}
Pain~: bleu-noir --- la farine est surtout de l'amidon. Fromage~: pas
d'autre changement de couleur que la teinte propre du diiode ---
\omterm{def:g10:chemistry-of-life:families}{protéines} et \omterm{def:g10:chemistry-of-life:families}{lipides}, pas d'amidon. Sucre~: aucun changement --- le
saccharose est un sucre simple, pas de l'amidon~; l'eau iodée ne
détecte que la longue chaîne d'amidon.
\end{solution}

\begin{solution}{exo:g10:chemistry-of-life:10}
La matière sèche est surtout organique~: des molécules
carbone--hydrogène qui se combinent à l'oxygène en libérant de
l'énergie, du dioxyde de carbone et de l'eau. La cendre, ce sont les
\omterm{def:g10:chemistry-of-life:mineral-organic}{sels minéraux}, déjà entièrement combinés à l'oxygène~; rien en eux ne
peut brûler.
\end{solution}

\begin{solution}{exo:g10:chemistry-of-life:11}
La liqueur de Fehling réagit avec un groupement chimique libre porté
par les molécules de glucose isolées~; dans l'amidon, les unités de
glucose sont liées les unes aux autres et ces groupements sont
enfermés dans la chaîne, si bien que le réactif ne trouve rien. Le
diiode, à l'inverse, se loge dans la longue chaîne enroulée de
l'amidon, que des molécules de glucose isolées ne forment pas. Chaque
test détecte une structure, non un élément.
\end{solution}

\begin{solution}{exo:g10:chemistry-of-life:12}
Le grain a absorbé massivement de l'eau. À 13~\%, les réserves sont
solides et les enzymes ne peuvent pas travailler~: la chimie du grain
est éteinte et il survit en dormance. L'eau dissout les réserves,
laisse agir les enzymes et porte les produits vers l'embryon en
croissance~: la chimie ne redémarre qu'une fois l'eau abondante.
\end{solution}

\begin{solution}{exo:g10:chemistry-of-life:13}
Sel~: $9 \times 0.5 = \qty{4.5}{g}$~; glucose~: $1 \times 0.5 =
\qty{0.5}{g}$. De l'eau pure diluerait le plasma~; l'eau entrerait
alors dans les globules rouges, qui gonfleraient et éclateraient.
\end{solution}

\begin{solution}{exo:g10:chemistry-of-life:14}
Glycogène~: $500 \times 17 = \qty{8500}{kJ}$. \omterm{def:g10:chemistry-of-life:families}{Lipides}~: $10000 \times 38 =
\qty{3.8e5}{kJ}$. Pour détenir \qty{3.8e5}{kJ} sous forme de glycogène~: $380000/17
\approx \qty{22.4}{kg}$ de glycogène, plus trois fois cette masse
d'eau, environ \qty{90}{kg} en tout --- plus que la personne. La
graisse stocke à peu près neuf fois plus d'énergie par kilogramme
porté, et c'est pourquoi les animaux qui doivent se déplacer stockent
leurs réserves sous forme de \omterm{def:g10:chemistry-of-life:families}{lipides}.
\end{solution}

\begin{solution}{exo:g10:chemistry-of-life:15}
Les atomes sont bien ordinaires, mais (1) les proportions ne le sont
pas~: la matière vivante concentre le carbone d'un facteur mille et
l'azote bien davantage, face à une croûte d'oxygène et de silicium~;
(2) les molécules non plus~: seuls les êtres vivants assemblent le
carbone en chaînes géantes --- \omterm{def:g10:chemistry-of-life:families}{protéines}, amidon, ADN --- dont la
taille n'a pas d'équivalent minéral~; (3) ces molécules portent une
information dans l'ordre de leurs unités, ce qu'aucun cristal ne fait.
La vie est chimiquement particulière par son organisation, non par ses
atomes.
\end{solution}

\begin{solution}{pb:g10:chemistry-of-life:1}
\textbf{1.} Eau $0.13 \times 45 = \qty{5.9}{mg}$~; masse sèche
\qty{39.2}{mg}.

\textbf{2.} Amidon $0.70 \times 39.2 \approx \qty{27.4}{mg}$~;
\omterm{def:g10:chemistry-of-life:families}{protéines} $0.12 \times 39.2 \approx \qty{4.7}{mg}$~; \omterm{def:g10:chemistry-of-life:families}{lipides} $0.02 \times 39.2
\approx \qty{0.8}{mg}$.

\textbf{3.} Eau iodée (bleu-noir) pour l'amidon~; biuret (violet) pour
les \omterm{def:g10:chemistry-of-life:families}{protéines}~; une tache translucide sur papier ou le Soudan III pour
les \omterm{def:g10:chemistry-of-life:families}{lipides}.

\textbf{4.} À 13~\%, les réserves sont solides et les enzymes
inactives~; l'eau doit dissoudre les réserves et laisser les réactions
se faire avant que l'embryon puisse croître. La germination commence
par l'absorption d'eau, et les 88~\% de la plantule sont le signe que
sa chimie est allumée.

\textbf{5.} Eau $0.13 \times 8 = \qty{1.04}{t}$~; matière sèche
\qty{6.96}{t}, soit environ \qty{7}{t}.

\textbf{6.} Masse $500 + 320 + 10 - 130 = \qty{700}{g}$. Eau du pain~:
venant de la farine $0.13 \times 500 = \qty{65}{g}$, plus \qty{320}{g},
moins \qty{130}{g}~: \qty{255}{g}, soit $255/700 \approx 36\%$.

\textbf{7.} Matière sèche de la farine $0.87 \times 500 = \qty{435}{g}$~;
amidon $0.70 \times 435 \approx \qty{305}{g}$~; \omterm{def:g10:chemistry-of-life:families}{protéines} $0.12 \times 435 \approx
\qty{52}{g}$ (\omterm{def:g10:chemistry-of-life:families}{lipides} environ \qty{9}{g}).

\textbf{8.} $305 \times 17 + 52 \times 17 + 9 \times 38 \approx 5185 +
884 + 342 \approx \qty{6400}{kJ}$.

\textbf{9.} $120/700 \approx 0.17$ du pain~: environ \qty{1100}{kJ},
soit 11~\% de \qty{10000}{kJ}.

\textbf{10.} Le sel est minéral~: il est déjà entièrement oxydé et n'a
pas de liaison carbone--hydrogène à brûler. Dans le corps, ses ions
fixent la salinité du sang et sont nécessaires aux signaux nerveux et
musculaires.

\textbf{11.} $0.60 \times 60 = \qty{36}{kg}$ d'eau.

\textbf{12.} \omterm{def:g10:chemistry-of-life:families}{Protéines} $0.17 \times 60 = \qty{10.2}{kg}$~; \omterm{def:g10:chemistry-of-life:families}{lipides} $0.16
\times 60 = \qty{9.6}{kg}$~; minéraux \qty{3.6}{kg}~; \omterm{def:g10:chemistry-of-life:families}{glucides}
\qty{0.6}{kg}.

\textbf{13.} $1.2/36 \approx 3.3\%$ de son eau. Le volume sanguin et
le contrôle de la température en dépendent tous deux~; une perte de
10~\% est dangereuse, et la boisson doit donc la rétablir en quelques
heures.

\textbf{14.} Plasma $0.55 \times 4.5 \approx \qty{2.5}{L}$~; eau $0.9
\times 2.5 \approx \qty{2.2}{kg}$, soit environ 6~\% de l'eau du corps.
Le reste est dans les cellules (environ les deux tiers de toute l'eau
du corps) et dans le liquide entre les cellules.

\textbf{15.} Un nourrisson échange chaque jour une fraction bien plus
grande de son eau (boisson, urine, peau) qu'un adulte~; quelques
heures de pertes sans apport lui retirent donc une fraction dangereuse
de son total.

\textbf{16.} \omterm{def:g10:chemistry-of-life:families}{Glucides} $0.40 \times 0.6 = \qty{0.24}{kg}$~; \omterm{def:g10:chemistry-of-life:families}{lipides}
$0.77 \times 9.6 \approx \qty{7.4}{kg}$~; \omterm{def:g10:chemistry-of-life:families}{protéines} $0.53 \times 10.2
\approx \qty{5.4}{kg}$. Total environ \qty{13}{kg} de carbone.

\textbf{17.} $13/60 \approx 22\%$, contre 18.5~\%. L'accord est bon
compte tenu des fractions arrondies par famille et de la variation de
la teneur en \omterm{def:g10:chemistry-of-life:families}{lipides} d'un individu à l'autre~; un corps plus sec donne
un chiffre plus bas.

\textbf{18.} $0.16 \times 10.2 \approx \qty{1.6}{kg}$ d'azote. Il
venait de l'azote de l'air, capté par les bactéries du sol, prélevé
par les plantes et mangé le long de la chaîne alimentaire.

\textbf{19.} $0.89 \times 36 \approx \qty{32}{kg}$ d'oxygène, soit
$32/60 \approx 53\%$ de la masse du corps. L'eau à elle seule rend
compte de l'essentiel des 65~\%~; le reste est l'oxygène contenu dans
les molécules organiques.

\textbf{20.} Une personne de \qty{60}{kg} porte environ \qty{13}{kg}
de carbone. Il a quitté l'air sous forme de dioxyde de carbone par la
photosynthèse des plantes, et il est entré dans le corps par
l'alimentation et la digestion le long de la chaîne alimentaire.
\end{solution}
